\begin{derivation}[协变光子模式的对易关系]
在 Feynman 规范 $\xi_{\rm gf}=1$ 中，先把
\[
\mathcal L_{\rm EM,1}
=-\frac1{4\mu_0}F_{\mu\nu}F^{\mu\nu}
-\frac1{2\mu_0}(\partial_\mu A^\mu)^2
\]
展开。利用
\[
F_{\mu\nu}F^{\mu\nu}
=2\bigl[(\partial_\mu A_\nu)(\partial^\mu A^\nu)
-(\partial_\mu A_\nu)(\partial^\nu A^\mu)\bigr]
\]
并对第二项分部积分，交叉项与规范固定项相消；忽略不影响 Euler--Lagrange 方程的全微分后，二次拉氏量变成
\[
\mathcal L_{\rm EM,1}
=-\frac1{2\mu_0}(\partial_\mu A_\nu)(\partial^\mu A^\nu).
\]
因此以 $A^\nu$ 为正则坐标时
\[
\Pi_\nu
\equiv\frac{\partial\mathcal L_{\rm EM,1}}
{\partial(\partial_t A^\nu)}
=-\frac1{\mu_0c^2}\partial_tA_\nu
=-\varepsilon_0\partial_tA_\nu.
\]
这里的下指标十分重要：Minkowski 度规已经进入 $A_\nu=\eta_{\nu\rho}A^\rho$，所以四个分量不会得到同号的正定振子度量。

令
\[
N_{\vb{\kappa}}
\equiv\sqrt{\frac{\hbar}{2\varepsilon_0\omega_{\vb{\kappa}}}}
\]
并把正文\eqnrefs{eq:covariant-photon-fourpol-r7} 写成
\[
A^\mu(x)=\sum_\lambda\int\frac{\dd^3\kappa}{(2\pi)^3}N_{\vb{\kappa}}
\left(\epsilon_\lambda^\mu a_\lambda \ee^{-\ii\kappa\cdot x}
+\epsilon_\lambda^{\mu *}a_\lambda^\dagger \ee^{+\ii\kappa\cdot x}\right).
\]
时间微分给出 $\mp\ii\omega_{\vb{\kappa}}$，所以
\[
\Pi_\nu(x)=\ii\varepsilon_0\sum_\lambda\int\frac{\dd^3\kappa}{(2\pi)^3}
\omega_{\vb{\kappa}}N_{\vb{\kappa}}
\left(\epsilon_{\lambda\nu}a_\lambda \ee^{-\ii\kappa\cdot x}
-\epsilon_{\lambda\nu}^*a_\lambda^\dagger \ee^{+\ii\kappa\cdot x}\right).
\]
暂不预设模代数，而写
\[
[a_\lambda(\vb{\kappa}),a_{\lambda'}^\dagger(\vb{\kappa}')]
=C_{\lambda\lambda'}(\vb{\kappa})
(2\pi)^3\delta^{(3)}(\vb{\kappa}-\vb{\kappa}').
\]
在等时 $x^0=y^0$ 计算 $[A^\mu(\vb{x}),\Pi_\nu(\vb{y})]$。只有 $a$--$a^\dagger$ 与 $a^\dagger$--$a$ 两类收缩；用 $N_{\vb{\kappa}}^2\varepsilon_0\omega_{\vb{\kappa}}=\hbar/2$ 后，两类正负频贡献恰好合成
\begin{align*}
[A^\mu(\vb{x}),\Pi_\nu(\vb{y})]
&=-\ii\hbar\int\frac{\dd^3\kappa}{(2\pi)^3}
 \ee^{\ii\vb{\kappa}\cdot(\vb{x}-\vb{y})}
 \sum_{\lambda\lambda'}C_{\lambda\lambda'}
 \epsilon_\lambda^\mu\epsilon_{\lambda'\nu}^*.
\end{align*}
正则量子化要求这等于
\[
\ii\hbar\delta^\mu{}_{\nu}\delta^{(3)}(\vb{x}-\vb{y})
=\ii\hbar\int\frac{\dd^3\kappa}{(2\pi)^3}
 \ee^{\ii\vb{\kappa}\cdot(\vb{x}-\vb{y})}\delta^\mu{}_{\nu}.
\]
由于四偏振基满足正文\eqnrefs{eq:fourpol-completeness-r7}，即
\[
\sum_{\lambda\lambda'}\eta^{\lambda\lambda'}
\epsilon_\lambda^\mu\epsilon_{\lambda'\nu}^*=\delta^\mu{}_{\nu},
\]
逐个 Fourier 模比较得到唯一解
\[
{C_{\lambda\lambda'}=-\eta_{\lambda\lambda'}}.
\]
于是
\[
[a_\lambda(\vb{\kappa}),a_{\lambda'}^\dagger(\vb{\kappa}')]
=-\eta_{\lambda\lambda'}(2\pi)^3\delta^{(3)}(\vb{\kappa}-\vb{\kappa}'),
\]
即正文\eqnrefs{eq:covariant-photon-mode-comm-r7}。特别地，在这个不定内积空间中，单个时间样激发的自内积为
\[
\langle0|a_0a_0^\dagger|0\rangle<0,
\]
所以负范数并非另加的假设，而是“显式 Lorentz 协变 + 四分量正则量子化 + Minkowski 度规”三者共同固定的代数结果。Gupta--Bleuler 条件随后才负责从这个不定内积空间中选出物理商空间。
\end{derivation}

\begin{derivation}[横向场的正则量子化与模式归一化]
Maxwell 拉氏量

\begin{equation*}
L_{\rm EM}=\frac{\varepsilon_0}{2}\int \dd^3 x\,(\vb{E}^2-c^2\vb{B}^2)
\end{equation*}
给 $\Pi_\phi=0$ 与 $\vb{\Pi}=-\varepsilon_0\vb{E}$。Legendre 变换并对 $\varepsilon_0\vb{E}\cdot\nabla\phi$ 分部积分后
\begin{equation*}
H_{\rm EM}=\int \dd^3 x\left[
\frac{\varepsilon_0}{2}(\vb{E}^2+c^2\vb{B}^2)
+\phi\,\nabla\cdot\vb{\Pi}\right].
\end{equation*}
所以 $\phi$ 是 Lagrange 乘子，变分给 $\nabla\cdot\vb{\Pi}=0$。结合 Coulomb 规范，$A$ 与 $\Pi$ 都只能取横向分量。

写 $P_L=\nabla\nabla/\nabla^2$，则 $P_L^2=P_L$，所以 $P_T=I-P_L$ 满足 $P_T^2=P_T$ 与 $\nabla\cdot P_T=0$。把未施加约束的正则括号 $[A_i,\Pi_j]=\ii\hbar\delta_{ij}\delta^{(3)}$ 两侧各投影一次：
\begin{align*}
[A_i^T(\vb{x}),\Pi_j^T(\vb{y})]
&=(P_T)_{ik}^{(x)}(P_T)_{j\ell}^{(y)}
[A_k(\vb{x}),\Pi_\ell(\vb{y})]\\
&=\ii\hbar(P_T)_{ij}^{(x)}\delta^{(3)}(\vb{x}-\vb{y}),
\end{align*}
这就是\eqnrefs{eq:transverse-comm-r2}。

对固定 $\vb{\kappa}$，$\{\vb{\epsilon}_1,\vb{\epsilon}_2,\hat{\vb{\kappa}}\}$ 构成三维正交归一基，因此单位算符的完备分解给出
\begin{equation*}
\sum_{\lambda=1}^2\epsilon_{\lambda i}\epsilon_{\lambda j}^*
=\delta_{ij}-\hat\kappa_i\hat\kappa_j.
\end{equation*}
再把每个横向 Fourier 系数视为谐振子，以 $[a,a^\dagger]=(2\pi)^3\delta^{(3)}$ 归一化，并要求谐振子哈密顿量为 $\hbar\omega(a^\dagger a+1/2)$，就固定了\eqnrefs{eq:maxwell-mode-r2} 中的 $\sqrt{\hbar/(2\varepsilon_0\omega)}$ 系数。
\end{derivation}
